\documentclass[aspectratio=169,11pt]{beamer}

% --- Theme ---
\usetheme{metropolis}
\usepackage{appendixnumberbeamer}

% --- Packages ---
\usepackage{fontspec}
\usepackage{minted}
\setminted{fontsize=\small,breaklines=true}
\usepackage{booktabs}
\usepackage{tikz}
\usepackage{fontawesome5}

% --- Colors ---
\definecolor{healthgreen}{RGB}{34,139,34}
\definecolor{alertred}{RGB}{220,53,69}
\setbeamercolor{alerted text}{fg=alertred}

% --- Metadata ---
\title{Module 7: Regression for health outcomes}
\subtitle{Data analysis with Python for health specialists}
\author{Ya\'{e} Ulrich Gaba}
\institute{AIRINA Labs}
\date{2026}

\begin{document}

% ============================================================
\maketitle

% ============================================================
\begin{frame}{From correlation to prediction}
\begin{itemize}
  \item Chapter 6: ``Is there a \emph{difference}?'' (hypothesis testing)
  \item Chapter 7: ``Can we \emph{predict} one variable from another?''
\end{itemize}

\vspace{6pt}
Regression is the \alert{workhorse of clinical research}:
\begin{itemize}
  \item Every epidemiological study
  \item Every clinical trial analysis
  \item Every risk prediction model
\end{itemize}

\vspace{6pt}
\textbf{Key phrase:} ``after adjusting for'' = multiple regression.
\end{frame}

% ============================================================
\begin{frame}[fragile]{Simple linear regression}
\textbf{Q:} How does systolic BP change with age?

\begin{minted}{python}
from scipy import stats

slope, intercept, r, p, se = stats.linregress(
    df["age"], df["systolic_bp"])
print(f"SBP = {intercept:.1f} + {slope:.2f} x Age")
print(f"R-squared: {r**2:.3f}")
\end{minted}

\vspace{4pt}
Slope = 0.7 means: on average, SBP increases by \textbf{0.7~mmHg per year of age}.

\begin{exampleblock}{Caution}
This does \emph{not} mean aging \emph{causes} higher BP. It means the two are linearly associated in this sample.
\end{exampleblock}
\end{frame}

% ============================================================
\begin{frame}[fragile]{Multiple linear regression}
\textbf{Q:} Predict SBP from age, BMI, and smoking \emph{simultaneously}.

\begin{minted}{python}
import statsmodels.api as sm

X = sm.add_constant(df[["age", "bmi", "smoking"]])
model = sm.OLS(df["systolic_bp"], X).fit()
print(model.summary())
\end{minted}

\vspace{4pt}
\textbf{Interpreting coefficients:}
\begin{itemize}
  \item Age coef $\approx 0.6$: +0.6 mmHg per year, \emph{after adjusting for BMI and smoking}
  \item BMI coef $\approx 0.8$: +0.8 mmHg per BMI unit, \emph{after adjusting for age and smoking}
  \item Smoking coef $\approx 5$: smokers have 5 mmHg higher SBP, \emph{after adjusting for age and BMI}
\end{itemize}
\end{frame}

% ============================================================
\begin{frame}[fragile]{Confounders: the coffee trap}
\begin{minted}{python}
# Unadjusted: coffee appears to predict heart disease
r, p = stats.pearsonr(coffee, heart_disease_risk)
# r = 0.35, p < 0.001 -- significant!

# Adjusted: coffee effect disappears when controlling for smoking
X = sm.add_constant(df[["coffee", "smoking"]])
model = sm.OLS(df["hd_risk"], X).fit()
# Coffee p = 0.72 (not significant)
# Smoking p < 0.001 (the real cause)
\end{minted}

\begin{alertblock}{Lesson}
Without adjusting for smoking, you would wrongly conclude coffee increases heart disease risk. \alert{Always think about confounders} in observational studies.
\end{alertblock}
\end{frame}

% ============================================================
\begin{frame}[fragile]{Checking regression assumptions}
\begin{minted}{python}
fig, axes = plt.subplots(1, 3, figsize=(15, 4))

# 1. Residuals vs fitted (linearity + homoscedasticity)
axes[0].scatter(model.fittedvalues, model.resid, alpha=0.4)
axes[0].axhline(0, color="red", linestyle="--")

# 2. Histogram of residuals (normality)
axes[1].hist(model.resid, bins=25, edgecolor="black")

# 3. Q-Q plot (normality)
stats.probplot(model.resid, dist="norm", plot=axes[2])
\end{minted}

\textbf{4 assumptions:} linearity, independence, homoscedasticity, normality of residuals.
\end{frame}

% ============================================================
\begin{frame}[fragile]{Logistic regression: binary outcomes}
\textbf{Q:} Predict diabetes (yes/no) from age, BMI, and glucose.

\begin{minted}{python}
# statsmodels (for odds ratios and p-values)
X = sm.add_constant(pima[["age", "bmi", "glucose"]])
logit = sm.Logit(pima["outcome"], X).fit()
print(logit.summary())

# Odds ratios with 95% CI
import numpy as np
or_df = pd.DataFrame({
    "OR": np.exp(logit.params),
    "CI_low": np.exp(logit.conf_int()[0]),
    "CI_high": np.exp(logit.conf_int()[1]),
    "p": logit.pvalues
}).round(4)
\end{minted}
\end{frame}

% ============================================================
\begin{frame}{Interpreting odds ratios}
\begin{table}
\begin{tabular}{lcl}
\toprule
\textbf{Variable} & \textbf{OR} & \textbf{Interpretation} \\
\midrule
Glucose & 1.04 & +1 mg/dL $\to$ 4\% higher odds of diabetes \\
BMI & 1.08 & +1 kg/m$^2$ $\to$ 8\% higher odds \\
Age & 1.02 & +1 year $\to$ 2\% higher odds \\
\bottomrule
\end{tabular}
\end{table}

\vspace{6pt}
\begin{itemize}
  \item OR $> 1$: higher risk
  \item OR $< 1$: lower risk (protective)
  \item OR $= 1$: no association
\end{itemize}

All \alert{holding other variables constant}.
\end{frame}

% ============================================================
\begin{frame}[fragile]{Confusion matrix: evaluating the model}
\begin{minted}{python}
from sklearn.metrics import confusion_matrix
cm = confusion_matrix(y_test, y_pred)
tn, fp, fn, tp = cm.ravel()

sensitivity = tp / (tp + fn)  # recall
specificity = tn / (tn + fp)
ppv = tp / (tp + fp)          # precision
npv = tn / (tn + fn)
\end{minted}

\begin{alertblock}{Clinical priority}
\textbf{Screening} test: prioritize \alert{high sensitivity} (don't miss cases).\\
\textbf{Confirmatory} test: prioritize \alert{high specificity} (don't give false alarms).
\end{alertblock}
\end{frame}

% ============================================================
\begin{frame}[fragile]{Survival analysis: Cox model}
\begin{minted}{python}
from lifelines import CoxPHFitter

cph = CoxPHFitter()
cph.fit(surv_df, duration_col="time",
        event_col="event", formula="treatment")
cph.print_summary()

hr = np.exp(cph.params_["treatment"])
print(f"Hazard ratio: {hr:.3f}")
# HR = 0.65 -> treatment reduces hazard by 35%
\end{minted}

\vspace{4pt}
Hazard ratios are the standard in oncology trials, cardiovascular studies, and any time-to-event endpoint.
\end{frame}

% ============================================================
\begin{frame}{What you can do after this module}
\begin{enumerate}
  \item Fit simple and multiple linear regression with clinical interpretation
  \item Identify and control for confounders
  \item Fit logistic regression and interpret odds ratios
  \item Evaluate classification with sensitivity, specificity, PPV, NPV
  \item Run survival analysis with Kaplan-Meier and Cox models
\end{enumerate}

\vspace{12pt}
\centering
\Large \alert{Next: Module 8 --- Machine learning for clinical prediction}
\end{frame}

% ============================================================
\begin{frame}{Exercises (Hour 3)}
\begin{enumerate}
  \item \textbf{Simple logistic:} Predict diabetes from glucose only; report OR
  \item \textbf{Multiple logistic:} Add age, BMI, BP, pedigree; report all ORs with 95\% CIs
  \item \textbf{Thresholds:} Try classification thresholds of 0.3, 0.5, 0.7; compare sensitivity/specificity
  \item \textbf{Mini-project:} Build a complete diabetes risk prediction notebook with confusion matrix and clinical interpretation
\end{enumerate}
\end{frame}

\end{document}
